\chapter{Einleitung}
\label{chap:einleitung}

Die zunehmende Automatisierung verteilter Cloud- und Edge-Systeme stellt hohe Anforderungen an die Zuverlässigkeit der zugrundeliegenden Bereitstellungsmechanismen. Dieses einleitende Kapitel legt die Motivation der Untersuchung dar, ordnet das Vorhaben in den aktuellen Stand der Forschung ein, leitet daraus die zentrale Forschungsfrage sowie die Zielsetzung her und gibt einen Überblick über den Aufbau der Arbeit.

\section{Motivation}
\label{sec:motivation}

Die Ansprüche an heutige Cloud-Native-Infrastrukturen steigen stetig, vor allem an betriebsnahe Eigenschaften wie Ausfallsicherheit, automatisierte Selbstheilung und deterministische Reproduzierbarkeit. Genau für diese Metriken bilden Container-Orchestrierungssysteme wie Kubernetes in Verbindung mit deklarativen Bereitstellungsparadigmen das technische Fundament.\autocite{burns_borg_2016,burns_kubernetes_2022} Wurden Anwendungs- und Infrastrukturänderungen in klassischen \gl{ci}-/\gl{cd}-Pipelines historisch überwiegend über imperative, von außen angestoßene Deployment-Skripte nach dem Push-Prinzip in Zielumgebungen übertragen, hat sich mit GitOps ein gegenläufiges Betriebsmodell etabliert. Hierbei fungiert ein versionsverwaltetes Git-Repository als einzige verlässliche Wahrheitsquelle (\textit{\gl{ssot}}), während spezialisierte Software-Agenten innerhalb des Clusters den realen Systemzustand kontinuierlich mit dem deklarierten Soll-Zustand abgleichen.\autocite{beetz_gitops_2022,opengitops_principles_2021}

Aufgrund der zunehmenden Verbreitung leichtgewichtiger Kubernetes-Distributionen wie K3s verlagert sich der Einsatz solcher GitOps-Control-Planes zunehmend aus hochverfügbaren Rechenzentren in ressourcenbeschränkte Edge-, Zweigstellen- und Cloud-Instanzen.\autocite{k3s_architecture_2026,boehm_k3s_2021,koziolek_k3s_2023} Dies zeigt sich sowohl beim Betrieb verteilter \gl{iot}- und Edge-Knoten als auch bei kostenoptimierten Single-Node- und Kleinst-Clustern, die über öffentliche \gl{wan}-Strecken an externe Versionsverwaltungssysteme angebunden sind. In diesen Topologien kann eine unterbrechungsfreie Netzwerkverbindung zwischen dem Cluster-internen GitOps-Controller und dem externen Git-Repository nicht als dauerhaft gegeben vorausgesetzt werden. Außerdem gewinnt die Frage nach der Ressourcenstabilität der Control Plane auf Systemen mit strikt limitiertem Hauptspeicher (\gl{ram}) immer mehr an Bedeutung, da wiederholte Synchronisationsversuche und Timeouts bei Verbindungsstörungen zu Lastspitzen führen können. Diese Verschiebung der Einsatzgebiete rückt das dynamische Verhalten von GitOps-Operatoren unter transienter Netzwerkdegradation in den Fokus der Forschung und Systemarchitektur.\autocite{fogli_k3s_2021,ghosh_gitops_2025,basiri_chaos_2016}

\section{Stand der Forschung und Abgrenzung}
\label{sec:stand-der-forschung}

Die wissenschaftliche Literatur untersucht die Leistungsfähigkeit und Fehlertoleranz Cloud-nativer Orchestrierungs- und Bereitstellungssysteme bislang überwiegend in drei getrennten Teilbereichen. Einen ersten Schwerpunkt bildet die Evaluierung leichtgewichtiger Kubernetes-Distributionen auf ressourcenbeschränkten Knoten. Böhm und Wirtz profilieren den Ressourcenbedarf von MicroK8s und K3s im Vergleich zu einer regulären Kubernetes-Installation entlang typischer Lebenszyklusphasen und weisen für K3s einen deutlich geringeren Haupt- und Massenspeicherbedarf der Basisschicht nach.\autocite{boehm_k3s_2021} Koziolek und Eskandani erweitern diesen Vergleich auf MicroK8s, K3s, k0s und MicroShift und quantifizieren sowohl die minimale Speicherauslastung im Leerlauf als auch den Durchsatz der Control Plane unter Last.\autocite{koziolek_k3s_2023} Beide Studien setzen jedoch durchgängig störungsfreie lokale Netzwerke voraus und betrachten keine zusätzlichen GitOps-Operatoren. Fogli et\,al. untersuchen Standard-Kubernetes, K3s und KubeEdge zwar gezielt unter Paketverlusten, erhöhten Latenzen und Bandbreitenbegrenzungen in taktischen Edge-Netzwerken, richten den Fokus jedoch ausschließlich auf die Cluster-interne Verbindung zwischen Control Plane und entfernten Worker-Knoten.\autocite{fogli_k3s_2021}

Ein zweiter Forschungsstrang befasst sich mit dem konzeptionellen und empirischen Vergleich deklarativer GitOps-Werkzeuge. Aufbauend auf der architektonischen Einordnung des Pull-Prinzips durch Beetz und Harrer legen Ghosh et\,al. einen reproduzierbaren Leistungsbenchmark für die Operatoren ArgoCD, Flux\,CD und ConfigSync im Kontext von \textit{Intent-Based Networking} vor.\autocite{beetz_gitops_2022,ghosh_gitops_2025} Die Autoren erfassen die Verarbeitungs- und Bereitstellungslatenz, den Determinismus der Synchronisation sowie den Prozessor- und Speicherbedarf in Einzel- und Mehrfachszenarien.\autocite{ghosh_gitops_2025} Sämtliche Messreihen erfolgen allerdings unter idealen Netzwerkbedingungen, sodass das Konvergenzverhalten bei transienten Verbindungsstörungen zum externen Git-Repository sowie die resultierenden Wiederherstellungszeiten (\gl{mttr}) unberücksichtigt bleiben.

Drittens adressieren neuere Arbeiten der Systemforschung die automatisierte Zuverlässigkeitsprüfung von Kubernetes-Controllern und Operatoren. Sun et\,al. entwickeln mit \textit{Sieve} ein Testverfahren, das gezielt verzögerte Zustandsansichten, verworfene Ereignisse und asynchrone Teilausfälle zwischen Controller und \texttt{kube-apiserver} injiziert, um Konsistenz- und Konvergenzfehler in Reconciliation-Schleifen offenzulegen.\autocite{sun_sieve_2022} Daran anknüpfend überprüft \textit{Acto} von Gu et\,al. systematisch Zustandsübergänge von Kubernetes-Operatoren auf fehlerhafte Endzustände.\autocite{gu_acto_2023} Beide Ansätze zielen jedoch auf die qualitative Identifikation logischer Implementierungsfehler an der Cluster-internen \gl{api}-Grenze ab und erfassen weder reale Transportschicht-Störungen auf ausgehenden \gl{wan}-Strecken noch hochaufgelöste Latenz- und Ressourcenzeitreihen.

In der Zusammenschau bleibt die Schnittstelle zwischen ressourcenbeschränkten K3s-Umgebungen, deklarativen GitOps-Control-Planes und transienter Netzwerkdegradation empirisch weitgehend offen. Insbesondere fehlt eine quantitative Untersuchung, die den \textit{ArgoCD Application Controller} und den \textit{Repository-Server} selbst als \gl{sut} isoliert und die Kausalkette von probabilistischen Paketverlusten auf ausgehenden \gl{https}-Verbindungen (\gl{tcp}-Port~443) über Kernel-Retransmissionen und \gl{grpc}-Timeouts bis hin zur Reconcile-Latenz und Post-Fault-\gl{mttr} vermisst. Ebenso bleibt in bestehenden Benchmarks ungeklärt, ob wiederholte Synchronisationsversuche und blockierte Netzwerkroutinen auf speicherlimitierten Knoten zu Auslagerungsprozessen auf den Sekundärspeicher führen, welche Latenzmessungen durch Paging-Artefakte verfälschen würden.\autocite{gregg_systems_2020}

\section{Forschungsfrage und Zielsetzung}
\label{sec:forschungsfrage-zielsetzung}

Aus dieser Forschungslücke ergibt sich die Frage nach der quantifizierbaren Widerstandsfähigkeit (Resilienz) und dem Konvergenzverhalten moderner GitOps-Controller unter realen Fehlerbedingungen. Die vorliegende Arbeit widmet sich daher der folgenden Forschungsfrage:

\begin{quote}
\textit{Wie wirken sich transiente Paketverluste auf der Netzwerkverbindung zum externen Git-Repository auf die Reconcile-Latenz, die Zustandsübergänge, die Mean Time to Recovery (\gl{mttr}) sowie die Ressourcenauslastung des ArgoCD Application Controllers in einer ressourcenbeschränkten K3s-Umgebung aus?}
\end{quote}

Bei der fundierten Analyse dieser Forschungsfrage wird die Untersuchung nicht ausschließlich auf der Ebene einer qualitativen Funktionsprüfung geführt. Die Bewertung erfolgt gezielt aus der Perspektive der geschlossenen Regelungsschleife (\textit{Reconciliation Loop}), der zeitlichen Verteilung der Synchronisationslatenzen (Median und 95.-Perzentil), der Verweildauer in fehlerhaften Systemzuständen sowie der resultierenden \gl{cpu}-, Speicher- und Kernel-Paging-Auslastung.\autocite{beyer_sre_2016,nygard_release_2018} Hierdurch soll objektiv und reproduzierbar untersucht werden, inwiefern der \textit{ArgoCD Application Controller} (Version~3.6) bei massiven transienten \gl{tcp}-Paketverlusten seine Selbstheilungsfähigkeit wahrt, welche Verzögerungen durch exponentielle Backoff-Mechanismen und Kernel-Retransmissionen entstehen und ob die Stabilität des Gesamtsystems ohne Auslagerungsprozesse auf den Sekundärspeicher (\textit{Swap-Reinheit}) gewährleistet bleibt.\autocite{argocd_architecture_2026,gregg_systems_2020}

\section{Aufbau der Arbeit}
\label{sec:aufbau-der-arbeit}

Zur Beantwortung der Forschungsfrage ist die Arbeit in sechs Kapitel strukturiert. In Kapitel~\ref{chap:theoretische-grundlagen} werden die theoretischen Grundlagen deklarativer Infrastrukturmodelle, die Prinzipien des GitOps-Paradigmas, die regelungstheoretische Funktionsweise des Kubernetes-Operator-Patterns sowie die methodischen Grundlagen des Chaos Engineerings und der Resilienzmetriken erläutert.

Kapitel~\ref{chap:architektur-k3s-argocd} analysiert anschließend die spezifische Systemarchitektur der untersuchten Plattformkomponenten. Der Fokus liegt hierbei auf der leichtgewichtigen Kubernetes-Distribution K3s sowie der internen Mikroarchitektur, Zustandsmaschine und Fehlerbehandlung von ArgoCD im Kontext von Netzwerkstörungen und strikten Speichergrenzen.

Das vierte Kapitel beschreibt die Konzeption und Implementierung des im Rahmen dieser Arbeit entwickelten, vollständig als Code definierten Messframeworks (\textit{Infra-Chaos-Lab}). Dabei wird dargelegt, wie die ephemere Testumgebung auf Basis von Terraform und \gl{aws} \gl{ec2} automatisiert bereitgestellt, mit K3s, ArgoCD und Prometheus instrumentiert und um eine deterministische Fehlerinjektion auf Linux-Kernel-Ebene (\texttt{iptables}) sowie eine automatisierte Prüfung der Swap-Reinheit ergänzt wurde. Das fünfte Kapitel präsentiert und diskutiert die empirischen Messergebnisse der durchgeführten Chaos-Experimente anhand hochaufgelöster Zeitreihen- und Verteilungsanalysen. Die Arbeit endet im sechsten Kapitel mit einem konkreten Fazit, das die Forschungsfrage auf Basis der empirischen Befunde beantwortet, methodische Grenzen der Untersuchung reflektiert und einen Ausblick auf zukünftige Entwicklungen gibt.
